\section{Diseño, producción y resultados complementarios}
\label{anx:resultados}
% Origen: 08-paso7, 09-paso8 y 10-paso9 (tablas secundarias, sensibilidad, días simulados, NAC y n*), movidos al anexo.

\subsection{Escenarios del experimento}
\label{anx:escenarios}

La tabla~\ref{tab:escenarios-experimento} detalla la configuración de los cuatro escenarios comparados. Cruzan
los dos factores que la cátedra y el caso dejan en manos del operador: cuántos molinetes se habilitan y qué tan
rápido valida cada pasajero.

\begin{table}[H]
\centering
\caption{Escenarios del experimento}
\label{tab:escenarios-experimento}
\small
\begin{tabularx}{\textwidth}{@{}l l l X@{}}
\toprule
Escenario & Molinetes & Validación (s) & Qué representa \\
\midrule
\EO{} & 20 & Triangular (1,8; 2,4; 3,5) & Situación actual: los molinetes que el dato muestra en uso \\
\EI{} & 22 & Triangular (1,8; 2,4; 3,5) & Habilitar los 22 molinetes del vestíbulo Principal \\
\ET{} & 20 & Triangular (1,5; 1,8; 2,5) & Validación sin contacto (EMV o QR) con los molinetes actuales [supuesto] \\
\EIT{} & 22 & Triangular (1,5; 1,8; 2,5) & Las dos medidas juntas \\
\bottomrule
\end{tabularx}
\end{table}

La validación rápida toma como moda el extremo optimista del rango de la tabla~\ref{tab:parametros}
\parencite[Exh.~10-27]{tcqsm2013}. No hay medición propia del molinete EMV/QR que ya existe en la estación,
así que \ET{} y \EIT{} se leen como ``qué pasaría si la validación bajara a ese rango''.

\subsection{Sensibilidad}
\label{anx:sensibilidad}

Fuera de la familia principal se corren los contextos optimista y pesimista de la
tabla~\ref{tab:parametros} y, con el horario oficial, un atraso uniforme de los trenes de 0 y 240~s, para
ver si la conclusión cambia de signo.

\subsection{Días simulados}
\label{anx:dias}

La tabla~\ref{tab:dias-simulados} lista los días que sortearon las seis semillas. Los trenes son los que
llegaron a Plaza Constitución entre las 07:00 y las 09:30 de ese día. Los pasajeros generados reproducen las
validaciones observadas ese día: la diferencia es la carga de Poisson de cada tren.

\begin{table}[H]
\centering
\caption{Días hábiles simulados en la producción}
\label{tab:dias-simulados}
\small
\begin{tabular}{@{}llrrrr@{}}
\toprule
Semilla & Día & Trenes & Validaciones SBASE & Generados & Atascados (\EO{}) \\
\midrule
\num{20260923} & mié 13/05/2026 & 49 & \num{20390} & \num{20530} & 49 \\
\num{20260924} & vie 05/06/2026 & 48 & \num{17459} & \num{17457} & 26 \\
\num{20260925} & mié 25/03/2026 & 50 & \num{19781} & \num{19706} & 50 \\
\num{20260926} & vie 06/03/2026 & 49 & \num{16666} & \num{16760} & 22 \\
\num{20260927} & vie 15/05/2026 & 43 & \num{17847} & \num{17841} & 52 \\
\num{20260928} & mié 11/03/2026 & 50 & \num{17791} & \num{17620} & 26 \\
\bottomrule
\end{tabular}
\end{table}

\subsection{Variación entre días y medidas secundarias}
\label{anx:variacion}

\begin{figure}[H]
\centering
\figuraSiExiste{tpi-subte-res-boxplot.png}{0.75\textwidth}
\caption{P90 de espera por escenario. Cada punto es una réplica (un día hábil); las líneas unen los cuatro
escenarios del mismo día.}
\label{fig:boxplot}
\end{figure}

\textbf{Variación entre días.} El P90 de \EO{} va de 14,7~s (11/03) a 29,7~s (25/03), con un desvío de 6,0~s
entre días. El día con más pasajeros, el 13/05 con \num{20530}, no es el de mayor espera. La espera depende de
cómo se agrupan los trenes y no solo del total del día, que es el argumento que justifica simular las
oleadas. El orden de los escenarios, en cambio, se repite los seis días: $\EIT{} < \EI{} \approx \ET{} < \EO{}$
(figura~\ref{fig:boxplot}).

\textbf{Secundarias.} Lq, cola máxima y espera media bajan en el mismo orden que las primarias. Dos medidas no
sirven para comparar escenarios:
\begin{itemize}
  \item El P90 del tránsito del Roca hasta la cola es de unos 765~s en los cuatro escenarios. Es el recorrido
        por el andén, el hall y la escalera, y no depende de los molinetes.
  \item El tiempo de drenaje final queda en múltiplos de 900~s (1800, 2700 o 3600~s): lo fija el cierre por
        atasco, que revisa cada 15~minutos. Es una cota inferior y no se interpreta.
\end{itemize}

\subsection{Efecto de los números aleatorios comunes}
\label{anx:nac}

Con números aleatorios comunes, un día cargado es malo para los cuatro escenarios, así que las medidas de una
misma réplica están correlacionadas positivamente y la varianza de la diferencia queda por debajo de la suma de
las varianzas. La tabla~\ref{tab:nac} lo cuantifica para el P90 de espera: la correlación entre \EO{} y cada
alternativa es de 0,98, y la varianza de la diferencia queda entre un 43 y un 76~\% por debajo de la que
tendrían corridas independientes. Con la misma varianza por escenario y corridas independientes, los intervalos
habrían sido entre 1,3 y 2 veces más anchos ($\sqrt{39{,}67/16{,}95} = 1{,}53$ para \EI{}).

\begin{table}[H]
\centering
\caption{Efecto de los números aleatorios comunes sobre el P90 de espera (6 días)}
\label{tab:nac}
\small
\begin{tabular}{@{}lrrrrr@{}}
\toprule
Alternativa & $\operatorname{Cov}(X_{\EO{}}, X_a)$ & Correlación & $\operatorname{Var}(D)$ &
  $\operatorname{Var}(X_{\EO{}}) + \operatorname{Var}(X_a)$ & Reducción \\
\midrule
\EI{}  & \num{11.36} & \num{0.981} & \num{16.95} & \num{39.67} & 57~\% \\
\ET{}  & \num{16.62} & \num{0.984} & \num{10.65} & \num{43.89} & 76~\% \\
\EIT{} & \num{8.14}  & \num{0.977} & \num{21.59} & \num{37.88} & 43~\% \\
\bottomrule
\end{tabular}

\smallskip
\footnotesize Varianzas en s$^2$, con $\operatorname{Var}(X_{\EO{}}) = \num{35.94}$. Reducción
$= 1 - \operatorname{Var}(D) / [\operatorname{Var}(X_{\EO{}}) + \operatorname{Var}(X_a)]$.
\end{table}

\subsection{Cantidad de réplicas necesaria}
\label{anx:n-necesario}

La tabla~\ref{tab:n-necesario} da el $n^{*}$ que haría falta para cada primaria y alternativa, con los
semianchos fijados antes de correr: 1~s para los percentiles y 0,01 para la proporción.

\entradaGenerada{tabla-n-necesario}

Con los semianchos fijados antes de correr, $n^{*}$ va de 87 a 338 días para los percentiles y de 20 a 54 para
la proporción. Son cifras inalcanzables: la variación entre días es grande frente a un semiancho de 1~s, y cada
tanda de 10 corridas lleva del orden de 100 minutos. Completar los 10 días del diseño tampoco alcanzaría
esa precisión.

Se decidió, entonces, no agregar días. La pregunta del estudio es si cada alternativa mejora, y eso lo responde
el signo de los intervalos, que ya excluyen al cero los nueve con 6 días. Lo que queda impreciso es el tamaño de
la mejora, y así se informa: con su intervalo y no solo con la media. Como el $n$ quedó por debajo del $n_0 = 10$
planificado, el resultado se apoya en un $t$ crítico más exigente (4,655 con 5 grados de libertad, contra 3,622
con 9).
